\documentclass [conference]{IEEEtran}

\usepackage[utf8]{inputenc}
\usepackage[LGR,T1]{fontenc}
\usepackage[greek,english]{babel}
\usepackage{textalpha}
\DeclareFontFamilySubstitution
{LGR}{ptm}{Tempora-TLF}
\DeclareFontShape{LGR}{ptm}{m}{sc}
{<->ssub * Tempora-TLF/m/n}{}

\usepackage{graphicx} % Required for inserting images
\usepackage{float}



\title{1η Προγραμματιστική Εργασία\\[0.3em]
       {\large Αναζήτηση εικόνων με διανυσματικές αναπαραστάσεις}}
\author{
    \IEEEauthorblockN{ΛΑΖΑΡΟΣ ΛΑΖΑΡΙΔΗΣ}
    \IEEEauthorblockA{Α.Μ.: 1115201900359}
    \and
    \IEEEauthorblockN{ΠΑΝΑΓΙΩΤΗΣ ΜΑΜΟΥΖΕΛΟΣ}
    \IEEEauthorblockA{Α.Μ.: 1115202100091}
}
\date{October 2026}

\begin{document}

\maketitle
\section{Μετρήσεις:}

\begin{align}
    Οι μετρήσεις που έγιναν για την σύγκριση των αλγορίθμων είναι:\newline
    \item -Recall@5 
    \item -Recall@10
    \newline (Aφορούν πόσες είκονες που έπρεπε να βρεθούν, όντως ήταν εντός τον πρώτων 5,10 αντίστοιχα)
    \item -mAP@10\
    \newline (Πόσες βρέθηκαν στις 10, με μεγαλύτερη βαρύτητα στην σειρά την οποία βρέθηκαν)
    \item -Μέσος χρόνος αναζήτησης
    \newline
    
    Έγινε έλεγχος των αποτελεσμάτων για βάσεις με διαφορετικό μέγεθος: D = \{1000, 5000, 10000, 25000\}.
    \newline
    Κατά την διάρκεια ελέγχου για τα διάφορα μεγέθοι βάσεων, το μέγεθος του λεξιλογίου παραμένει σταθερό στo Κ = 256 και ο αριθμός των Sift Descriptors επίσης παραμένει S = 500. 
    \newline

    Στις συναρτήσεις LSH, Hypercube, IVF FLAT και IVF PQ, οι τιμές ορισμάτων τους που χρησημοποιούνται είναι οι προκαθορισμένες. \newline

    Έπειτα, για μέγεθος D = 10000, εξετάσαμε πως οι αλλαγές στο μέγεθος του λεξιλογίου, Κ = 512 καθώς και του αριθμού των τυχαίων Sift Descriptors που επιλέγονται, S = 1000 \newline

    Το seed των τυχαίων διαδικασιών παραμένει σταθερό σε όλους του υπολογισμούς.\newline

    Όσον αφορά τα iterations στην εκπαίδευση του kmeans:
\begin{itemize}
    \itemΓια K = 128 και S=100: \newline
    Yπήρχαν συνολικά 42000 σημεία. Στο iteration 50 άλλαξαν κέντρο 81 σημεία, περίπου 0.19\%
    \newline
    \itemΓια Κ = 512 και S = 100:\newline
    Yπήρχαν συνολικά 42000 σημεία. Σταμάτησε ο αλγόριθμος στο iteration 47 επειδή δεν άλλαξε κανένα σημείο κέντρο.
    \newline
    \itemΓια Κ = 512 και S = 1000: 
    Yπήρχαν 396255 σημεία. Στο iteration 50 άλλαξαν κέντρο 1703 σημεία, περίπου 0.42\%
    \newline
    \itemΓια Κ = 512 και S = 250: \newline
    Yπήρχαν 104822 σημεία. Στο iteration 50 άλλαξαν κέντρο 60 σημεία, περίπου 0.05\%
    \newline
\end{itemize}
    Ουσιαστικά βλέπουμε ότι μετά το iteration 20, πολύ λίγα σημεία αλλάζουν κέντρο και έτσι δεν θα είχε διαφορετικό αποτέλεσμα να χρησιμοποιούμε πάνω από 50 iterations.
    
\end{align}

\section{Αποτελέσματα για μέγεθος βάσης 1000}

\begin{itemize}

    \item Exact Search αποτελέσματα:
    \begin{table}[H]
    \centering
    \small
    \setlength{\tabcolsep}{3pt}
    \begin{tabular}{l r r r r r}
    \hline
    \textbf{Group} & \textbf{Queries} & \textbf{Recall@5} & \textbf{Recall@10} & \textbf{mAP@10} & \textbf{Time(ms)} \\
    \hline
    all & 23 & 0.7739 & 0.8087 & 0.7853 & 0.2124 \\
    illumination & 11 & 0.7273 & 0.7818 & 0.7511 & 0.2124 \\
    viewpoint & 12 & 0.8167 & 0.8333 & 0.8167 & 0.2125 \\
    \hline
    \end{tabular}
    \par\smallskip
    \footnotesize Queries με λιγότερα από 10 αποτελέσματα: 0
    \end{table}
    
    % \includegraphics[width=1\linewidth]{exact.256.500/exact1k.png}

    \item LSH αποτελέσματα:
    \begin{table}[H]
    \centering
    \small
    \setlength{\tabcolsep}{3pt}
    \begin{tabular}{l r r r r r}
    \hline
    \textbf{Group} & \textbf{Queries} & \textbf{Recall@5} & \textbf{Recall@10} & \textbf{mAP@10} & \textbf{Time(ms)} \\
    \hline
    all & 23 & 0.4261 & 0.4348 & 0.4278 & 0.0294 \\
    illumination & 11 & 0.2545 & 0.2545 & 0.2545 & 0.0276 \\
    viewpoint & 12 & 0.5833 & 0.6000 & 0.5867 & 0.0311 \\
    \hline
    \end{tabular}
    \par\smallskip
    \footnotesize Queries με λιγότερα από 10 αποτελέσματα: 6
    \end{table}

    % \includegraphics[width=1\linewidth]{lsh.256.500/lsh1k.png}

    \item Hypercube αποτελέσματα:
    \begin{table}[H]
    \centering
    \small
    \setlength{\tabcolsep}{3pt}
    kproj = 10, w = 0.05, M = 500, probes = 10
    \par\smallskip
    \begin{tabular}{l r r r r r}
    \hline
    \textbf{Group} & \textbf{Queries} & \textbf{Recall@5} & \textbf{Recall@10} & \textbf{mAP@10} & \textbf{Time(ms)} \\
    \hline
    all & 23 & 0.2870 & 0.3130 & 0.2860 & 0.0167 \\
    illumination & 11 & 0.2000 & 0.2364 & 0.1981 & 0.0148 \\
    viewpoint & 12 & 0.3667 & 0.3833 & 0.3667 & 0.0183 \\
    \hline
    \end{tabular}
    \par\smallskip
    \footnotesize Queries με λιγότερα από 10 αποτελέσματα: 2
    \end{table}

    % \includegraphics[width=1\linewidth]{hypercube.256.500/hypercube1k.png}

    \item IVF FLAT αποτελέσματα:
    \begin{table}[H]
    \centering
    \small
    \setlength{\tabcolsep}{3pt}
    34 clusters, nprobe = 5
    \par\smallskip
    \begin{tabular}{l r r r r r}
    \hline
    \textbf{Group} & \textbf{Queries} & \textbf{Recall@5} & \textbf{Recall@10} & \textbf{mAP@10} & \textbf{Time(ms)} \\
    \hline
    all & 23 & 0.7739 & 0.8087 & 0.7853 & 0.1488 \\
    illumination & 11 & 0.7273 & 0.7818 & 0.7511 & 0.1480 \\
    viewpoint & 12 & 0.8167 & 0.8333 & 0.8167 & 0.1495 \\
    \hline
    \end{tabular}
    \par\smallskip
    \footnotesize Queries με λιγότερα από 10 αποτελέσματα: 0
    \end{table}

    % \includegraphics[width=1\linewidth]{ivf.256.500/ivf1k.png}
        
    \item IVF PQ αποτελέσματα:
        \begin{table}[H]
        \centering
        \small
        \setlength{\tabcolsep}{3pt}
        34 clusters, nprobe = 5, M = 16
        \par\smallskip
        \begin{tabular}{l r r r r r}
        \hline
        \textbf{Group} & \textbf{Queries} & \textbf{Recall@5} & \textbf{Recall@10} & \textbf{mAP@10} & \textbf{Time(ms)} \\
        \hline
        all & 23 & 0.6957 & 0.7826 & 0.7334 & 0.1857 \\
        illumination & 11 & 0.6545 & 0.7636 & 0.7082 & 0.1884 \\
        viewpoint & 12 & 0.7333 & 0.8000 & 0.7565 & 0.1832 \\
        \hline
        \end{tabular}
        \par\smallskip
        \footnotesize Queries με λιγότερα από 10 αποτελέσματα: 0
        \end{table}

    % \includegraphics[width=1\linewidth]{ivfpq.256.500/ivfpq1k.png}
    
\end{itemize}

Για μέγεθος βάσης 1000 βλέπουμε πως ο αλγόριθμος exact search και ο IVF FLAT έχουν την μεγαλύτερη ακρίβεια. Ο IVF FLAT όμως έχει μέσο χρόνο εκτέλεσης 0.1488ms έναντι του 0.2124ms του exact search. Οι 2 ταχύτεροι αλγόριθμοι, Hypercube και LSH, με μέσο όρο 0.0167ms και 0.0294ms αντίστοιχα, παρουσιάζουν σημαντικά μειωμένες αποδόσεις στο Recall@n και mAP@10.O IFV PQ φένεται να είναι ο πιο ισοροπημένος καθώς έχει μετρήσεις που βρίσκονται ανάμεσα στα 2 άκρα των άλλων αλγορίθμων. Επιπλέον, παρατηρείται πως στους 2 ταχύτερους αλγορίθμους(Hypercube/LSH), η ακρίβεια τους στο σύνολο, είναι μειωμένη κυρίως λόγω του συνόλου εικόνων που είναι οι παραλλαγές φωτισμού, το οποίο απέχει σημαντικά από το μέσο όρο τους στο γενικό σύνολο.

\section{Αποτελέσματα για μέγεθος βάσης 5000}

\begin{itemize}

    \item Exact Search αποτελέσματα:
    \begin{table}[H]
    \centering
    \small
    \setlength{\tabcolsep}{3pt}
    \begin{tabular}{l r r r r r}
    \hline
    \textbf{Group} & \textbf{Queries} & \textbf{Recall@5} & \textbf{Recall@10} & \textbf{mAP@10} & \textbf{Time(ms)} \\
    \hline
    all & 23 & 0.7391 & 0.8000 & 0.7639 & 0.9721 \\
    illumination & 11 & 0.6545 & 0.7818 & 0.7124 & 0.9740 \\
    viewpoint & 12 & 0.8167 & 0.8167 & 0.8111 & 0.9703 \\
    \hline
    \end{tabular}
    \par\smallskip
    \footnotesize Queries με λιγότερα από 10 αποτελέσματα: 0
    \end{table}
    
    % \includegraphics[width=1\linewidth]{exact.256.500/exact5k.png}

    \item LSH αποτελέσματα:
    \begin{table}[H]
    \centering
    \small
    \setlength{\tabcolsep}{3pt}
    \begin{tabular}{l r r r r r}
    \hline
    \textbf{Group} & \textbf{Queries} & \textbf{Recall@5} & \textbf{Recall@10} & \textbf{mAP@10} & \textbf{Time(ms)} \\
    \hline
    all & 23 & 0.4261 & 0.4261 & 0.4261 & 0.0473 \\
    illumination & 11 & 0.2545 & 0.2545 & 0.2545 & 0.0457 \\
    viewpoint & 12 & 0.5833 & 0.5833 & 0.5833 & 0.0489 \\
    \hline
    \end{tabular}
    \par\smallskip
    \footnotesize Queries με λιγότερα από 10 αποτελέσματα: 2
    \end{table}

    % \includegraphics[width=1\linewidth]{lsh.256.500/lsh5k.png}

    \item Hypercube αποτελέσματα:
    \begin{table}[H]
    \centering
    \small
    \setlength{\tabcolsep}{3pt}
    kproj = 10, w = 0.05, M = 500, probes = 10
    \par\smallskip
    \begin{tabular}{l r r r r r}
    \hline
    \textbf{Group} & \textbf{Queries} & \textbf{Recall@5} & \textbf{Recall@10} & \textbf{mAP@10} & \textbf{Time(ms)} \\
    \hline
    all & 23 & 0.2783 & 0.2870 & 0.2739 & 0.0636 \\
    illumination & 11 & 0.2000 & 0.2000 & 0.1909 & 0.0548 \\
    viewpoint & 12 & 0.3500 & 0.3667 & 0.3500 & 0.0716 \\
    \hline
    \end{tabular}
    \par\smallskip
    \footnotesize Queries με λιγότερα από 10 αποτελέσματα: 0
    \end{table}

    % \includegraphics[width=1\linewidth]{hypercube.256.500/hypercube5k.png}

    \item IVF FLAT αποτελέσματα:
    \begin{table}[H]
    \centering
    \small
    \setlength{\tabcolsep}{3pt}
    72 clusters, nprobe = 5
    \par\smallskip
    \begin{tabular}{l r r r r r}
    \hline
    \textbf{Group} & \textbf{Queries} & \textbf{Recall@5} & \textbf{Recall@10} & \textbf{mAP@10} & \textbf{Time(ms)} \\
    \hline
    all & 23 & 0.7304 & 0.8000 & 0.7574 & 0.2679 \\
    illumination & 11 & 0.6545 & 0.7818 & 0.7133 & 0.2515 \\
    viewpoint & 12 & 0.8000 & 0.8167 & 0.7978 & 0.2830 \\
    \hline
    \end{tabular}
    \par\smallskip
    \footnotesize Queries με λιγότερα από 10 αποτελέσματα: 0
    \end{table}

    % \includegraphics[width=1\linewidth]{ivf.256.500/ivf5k.png}
        
    \item IVF PQ αποτελέσματα:
    \begin{table}[H]
    \centering
    \small
    \setlength{\tabcolsep}{3pt}
    72 clusters, nprobe = 5, Μ = 16
    \par\smallskip
    \begin{tabular}{l r r r r r}
    \hline
    \textbf{Group} & \textbf{Queries} & \textbf{Recall@5} & \textbf{Recall@10} & \textbf{mAP@10} & \textbf{Time(ms)} \\
    \hline
    all & 23 & 0.7043 & 0.7739 & 0.7243 & 0.2068 \\
    illumination & 11 & 0.6727 & 0.7455 & 0.6944 & 0.2024 \\
    viewpoint & 12 & 0.7333 & 0.8000 & 0.7517 & 0.2108 \\
    \hline
    \end{tabular}
    \par\smallskip
    \footnotesize Queries με λιγότερα από 10 αποτελέσματα: 0
    \end{table}

    % \includegraphics[width=1\linewidth]{ivfpq.256.500/ivfpq5k.png}
    
\end{itemize}

Όταν το μέγεθος της βάσης ανεβαίνει στα 5000, ο αλγόριθμος LSH και ο hypercube διατηρούν παρόμοια αποτελέσματα, με τον μέσο χρόνο εκτέλεσης να είναι η μόνη αξιοσημείωτη αλλαγή, από τα 0.0294ms στα 0.0473ms και από 0.0167ms στα 0.0636ms αντίστοιχα. Παραμένουν όμως οι 2 με την μικρότερη ακρίβεια. Η μεγαλύτερη ποσοστιαία μεταβολή στον χρόνο, είναι του exact search όπου η αύξηση του μεγέθους, οδήγησε σε σχεδόν ισοποσοστιαία αύξηση του χρόνου εκτέλεσης. Συνεχίζει όμως να παραμένει ο αλγόριθμος με την μεγαλύτερη ακρίβεια με κοντινό 2ο τον IVF FLAT ο οποίως όμως έχει μέσο όρο αναζήτησης 0.2679ms έναντι των 0.9721ms.

\section{Αποτελέσματα για μέγεθος βάσης 10000}

\begin{itemize}

    \item Exact Search αποτελέσματα:
    \begin{table}[H]
    \centering
    \small
    \setlength{\tabcolsep}{3pt}
    \begin{tabular}{l r r r r r}
    \hline
    \textbf{Group} & \textbf{Queries} & \textbf{Recall@5} & \textbf{Recall@10} & \textbf{mAP@10} & \textbf{Time(ms)} \\
    \hline
    all & 23 & 0.7391 & 0.7913 & 0.7574 & 1.9543 \\
    illumination & 11 & 0.6545 & 0.7636 & 0.6987 & 1.9576 \\
    viewpoint & 12 & 0.8167 & 0.8167 & 0.8111 & 1.9512 \\
    \hline
    \end{tabular}
    \par\smallskip
    \footnotesize Queries με λιγότερα από 10 αποτελέσματα: 0
    \end{table}
    
    % \includegraphics[width=1\linewidth]{exact.256.500/exact10k.png}

    \item LSH αποτελέσματα:
    \begin{table}[H]
    \centering
    \small
    \setlength{\tabcolsep}{3pt}
    \begin{tabular}{l r r r r r}
    \hline
    \textbf{Group} & \textbf{Queries} & \textbf{Recall@5} & \textbf{Recall@10} & \textbf{mAP@10} & \textbf{Time(ms)} \\
    \hline
    all & 23 & 0.4261 & 0.4261 & 0.4261 & 0.0726 \\
    illumination & 11 & 0.2545 & 0.2545 & 0.2545 & 0.0659 \\
    viewpoint & 12 & 0.5833 & 0.5833 & 0.5833 & 0.0788 \\
    \hline
    \end{tabular}
    \par\smallskip
    \footnotesize Queries με λιγότερα από 10 αποτελέσματα: 0
    \end{table}

    % \includegraphics[width=1\linewidth]{lsh.256.500/lsh10k.png}

    \item Hypercube αποτελέσματα:
    \begin{table}[H]
    \centering
    \small
    \setlength{\tabcolsep}{3pt}
    kproj = 10, w = 0.05, M = 500, probes = 10
    \par\smallskip
    \begin{tabular}{l r r r r r}
    \hline
    \textbf{Group} & \textbf{Queries} & \textbf{Recall@5} & \textbf{Recall@10} & \textbf{mAP@10} & \textbf{Time(ms)} \\
    \hline
    all & 23 & 0.2087 & 0.2174 & 0.2041 & 0.0850 \\
    illumination & 11 & 0.1636 & 0.1636 & 0.1545 & 0.0739 \\
    viewpoint & 12 & 0.2500 & 0.2667 & 0.2494 & 0.0951 \\
    \hline
    \end{tabular}
    \par\smallskip
    \footnotesize Queries με λιγότερα από 10 αποτελέσματα: 0
    \end{table}

    % \includegraphics[width=1\linewidth]{hypercube.256.500/hypercube10k.png}

    \item IVF FLAT αποτελέσματα:
    \begin{table}[H]
    \centering
    \small
    \setlength{\tabcolsep}{3pt}
    101 clusters, nprobe = 5
    \par\smallskip
    \begin{tabular}{l r r r r r}
    \hline
    \textbf{Group} & \textbf{Queries} & \textbf{Recall@5} & \textbf{Recall@10} & \textbf{mAP@10} & \textbf{Time(ms)} \\
    \hline
    all & 23 & 0.6870 & 0.7304 & 0.7070 & 0.4334 \\
    illumination & 11 & 0.6182 & 0.7091 & 0.6662 & 0.4073 \\
    viewpoint & 12 & 0.7500 & 0.7500 & 0.7444 & 0.4574 \\
    \hline
    \end{tabular}
    \par\smallskip
    \footnotesize Queries με λιγότερα από 10 αποτελέσματα: 0
    \end{table}

    % \includegraphics[width=1\linewidth]{ivf.256.500/ivf10k.png}
        
    \item IVF PQ αποτελέσματα:
    \begin{table}[H]
    \centering
    \small
    \setlength{\tabcolsep}{3pt}
    101 clusters, nprobe = 5, M = 16
    \par\smallskip
    \begin{tabular}{l r r r r r}
    \hline
    \textbf{Group} & \textbf{Queries} & \textbf{Recall@5} & \textbf{Recall@10} & \textbf{mAP@10} & \textbf{Time(ms)} \\
    \hline
    all & 23 & 0.6435 & 0.6957 & 0.6445 & 0.2251 \\
    illumination & 11 & 0.6182 & 0.6364 & 0.5969 & 0.2225 \\
    viewpoint & 12 & 0.6667 & 0.7500 & 0.6882 & 0.2275 \\
    \hline
    \end{tabular}
    \par\smallskip
    \footnotesize Queries με λιγότερα από 10 αποτελέσματα: 0
    \end{table}

    % \includegraphics[width=1\linewidth]{ivfpq.256.500/ivfpq10k.png}
    
\end{itemize}

Στα αποτελέσματα για μέγεθος 10000 παρατηρούμε πως o LSH παραμένει σταθερός στην ακρίβεια του, ενώ ο hypercube κυρίως, στην ομάδα εικόνων που αφορούν την οπτική γωνία, παρουσιάζει μεγάλη μείωση της ακρίβειας του. Η αύξηση του μέσου χρόνου αναζήτησης του exact search, συνεχίζει να μεταβάλεται με όμοιo ρυθμό του μεγέθους της βάσης εικόνων και ο χρόνος του ξεπερνάει τους υπόλοιπους αλγορίθμους κατά ποσοστό τουλάχιστον 300%, καθώς ο 2ος στον μέσο χρόνο αναζήτησης είναι ο IFV FLAT ο οποίως διατηρεί ανταγωνιστικό χρόνο αναζήτησης ακόμη, στα 0.4334ms, αντί του 1.9543ms του exact search. Διατηρεί όμως ο exact search τις καλύτερες βαθμολογίες στην ακρίβεια του.


\section{Αποτελέσματα για μέγεθος βάσης 25000}

\begin{itemize}

    \item Exact Search αποτελέσματα:
    \begin{table}[H]
    \centering
    \small
    \setlength{\tabcolsep}{3pt}
    \begin{tabular}{l r r r r r}
    \hline
    \textbf{Group} & \textbf{Queries} & \textbf{Recall@5} & \textbf{Recall@10} & \textbf{mAP@10} & \textbf{Time(ms)} \\
    \hline
    all & 23 & 0.7304 & 0.7478 & 0.7230 & 4.7523 \\
    illumination & 11 & 0.6364 & 0.6727 & 0.6482 & 4.7557 \\
    viewpoint & 12 & 0.8167 & 0.8167 & 0.7917 & 4.7492 \\
    \hline
    \end{tabular}
    \par\smallskip
    \footnotesize Queries με λιγότερα από 10 αποτελέσματα: 0
    \end{table}
    
    % \includegraphics[width=1\linewidth]{exact.256.500/exact25k.png}

    \item LSH αποτελέσματα:
    \begin{table}[H]
    \centering
    \small
    \setlength{\tabcolsep}{3pt}
    \begin{tabular}{l r r r r r}
    \hline
    \textbf{Group} & \textbf{Queries} & \textbf{Recall@5} & \textbf{Recall@10} & \textbf{mAP@10} & \textbf{Time(ms)} \\
    \hline
    all & 23 & 0.4261 & 0.4261 & 0.4261 & 0.1550 \\
    illumination & 11 & 0.2545 & 0.2545 & 0.2545 & 0.1404 \\
    viewpoint & 12 & 0.5833 & 0.5833 & 0.5833 & 0.1685 \\
    \hline
    \end{tabular}
    \par\smallskip
    \footnotesize Queries με λιγότερα από 10 αποτελέσματα: 0
    \end{table}

    % \includegraphics[width=1\linewidth]{lsh.256.500/lsh25k.png}

    \item Hypercube αποτελέσματα:
    \begin{table}[H]
    \centering
    \small
    \setlength{\tabcolsep}{3pt}
    kproj = 10, w = 0.05, M = 500, probes = 10
    \par\smallskip
    \begin{tabular}{l r r r r r}
    \hline
    \textbf{Group} & \textbf{Queries} & \textbf{Recall@5} & \textbf{Recall@10} & \textbf{mAP@10} & \textbf{Time(ms)} \\
    \hline
    all & 23 & 0.1826 & 0.1826 & 0.1754 & 01117 \\
    illumination & 11 & 0.1636 & 0.1636 & 0.1545 & 0.1065 \\
    viewpoint & 12 & 0.2000 & 0.2000 & 0.1944 & 0.1165 \\
    \hline
    \end{tabular}
    \par\smallskip
    \footnotesize Queries με λιγότερα από 10 αποτελέσματα: 0
    \end{table}

    % \includegraphics[width=1\linewidth]{hypercube.256.500/hypercube25k.png}

    \item IVF FLAT αποτελέσματα:
    \begin{table}[H]
    \centering
    \small
    \setlength{\tabcolsep}{3pt}
    159 clusters, nprobe = 5
    \par\smallskip
    \begin{tabular}{l r r r r r}
    \hline
    \textbf{Group} & \textbf{Queries} & \textbf{Recall@5} & \textbf{Recall@10} & \textbf{mAP@10} & \textbf{Time(ms)} \\
    \hline
    all & 23 & 0.6435 & 0.6435 & 0.6283 & 0.7662 \\
    illumination & 11 & 0.5273 & 0.5273 & 0.5227 & 0.7675 \\
    viewpoint & 12 & 0.7500 & 0.7500 & 0.7250 & 0.7650 \\
    \hline
    \end{tabular}
    \par\smallskip
    \footnotesize Queries με λιγότερα από 10 αποτελέσματα: 0
    \end{table}

    % \includegraphics[width=1\linewidth]{ivf.256.500/ivf25k.png}
        
    \item IVF PQ αποτελέσματα:
    \begin{table}[H]
    \centering
    \small
    \setlength{\tabcolsep}{3pt}
    159 clusters, nprobe = 5, M = 16
    \par\smallskip
    \begin{tabular}{l r r r r r}
    \hline
    \textbf{Group} & \textbf{Queries} & \textbf{Recall@5} & \textbf{Recall@10} & \textbf{mAP@10} & \textbf{Time(ms)} \\
    \hline
    all & 23 & 0.6174 & 0.6261 & 0.6108 & 0.2943 \\
    illumination & 11 & 0.5273 & 0.5455 & 0.5134 & 0.3156 \\
    viewpoint & 12 & 0.7000 & 0.7000 & 0.7000 & 0.2748 \\
    \hline
    \end{tabular}
    \par\smallskip
    \footnotesize Queries με λιγότερα από 10 αποτελέσματα: 0
    \end{table}

    % \includegraphics[width=1\linewidth]{ivfpq.256.500/ivfpq25k.png}
    
\end{itemize}

Τέλος, αυξάνοντας για μία ακόμη φορά το μέγεθος, πλέον στις 25000 εικόνες, βλέπουμε το κύριο πρόβλημα των exact search/IFV FLAT, στους οποίους ο μέσος χρόνος εκτέλεσης έχει πλέον αποκτήσει μεγάλο χάσμα συγκριτικά με τους υπόλοιπους. Ο IFV PQ παραμένει ο αλγόριθμος με τις μικρότερες μεταβολές στις μετρήσεις του ανεξαρτήτως του μεγέθους της βάσης. 

\section{Συμπεράσματα αποτελεσμάτων καθώς αλλάζει το μέγεθος της βάσης εικόνων}
\begin{itemize}
    

    \item Exact Search:\newline
    Ο αλγόριθμος υψηλότερης ακρίβειας, με το μεγαλύτερο χρονικό κόστος όμως. Καλύτερη επιλογή για βάσεις περιορισμένου μεγέθους όπου η ακρίβεια είναι σημαντική.

    \item LSH:\newline
    Από τους ταχύτερους μέσους όρους χρόνου αναζήτησης, με αρκετά μειωμένη ακρίβεια. Καλή επιλογή για εκτενείς βάσεις στις οποίες όμως η ακρίβεια της αναζήτησης δεν είναι τόσο σημαντική. Παρόλα αυτά, αν αλλάξουμε τις παραμέτρους που λαμβάνει για τo ποσό πινάκων/συναρτήσεων κατακερματισμού και το μέγεθος των κουβάδων μπορεί να αποκτήσει ικανοποιητική ακρίβεια, εις βάρος όμως της ταχύτητας του.

    \item Hypercube:\newline
    Ο ταχύτερος αλγόριθμος, με την μικρότερη ακρίβεια. Καλός για υπέρογκες βάσεις στις οποίες η ταχύτητα μετράει περισσότερο από την ακρίβεια.

    \item IVF FLAT:\newline
    Διατηρεί καλές επιδόσεις, όμοιες σχετικά του exact search, στην ακρίβεια, χωρίς το χρονικό του κόστος. Καλός για μικρομεσαίες βάσεις.

    \item IVF PQ: \newline
    Ο καλύτερος μέσος αλγόριθμος. Ικανοποιητική απόδοση σε όλα τα μεγέθοι βάσεων. Ιδανικός για βάσεις οι οποίες αλλάζουν μέγεθος συνεχώς.

\end{itemize}

\section{Επιρροή του λεξιλογίου}

\begin{itemize}
    \item Exact Search K = 128 and K = 512:
    \begin{table}[H]
    \centering
    \small
    \setlength{\tabcolsep}{3pt}
    \begin{tabular}{l r r r r r}
    \hline
    \textbf{Group} & \textbf{Queries} & \textbf{Recall@5} & \textbf{Recall@10} & \textbf{mAP@10} & \textbf{Time(ms)} \\
    \hline
    all & 23 & 0.7217 & 0.7565 & 0.7283 & 0.9824 \\
    illumination & 11 & 0.6545 & 0.6909 & 0.6739 & 0.9891 \\
    viewpoint & 12 & 0.7833 & 0.8167 & 0.7781 & 0.9763 \\
    \hline
    \end{tabular}
    \par\smallskip
    \footnotesize Queries με λιγότερα από 10 αποτελέσματα: 0
    \end{table}
    
    % \includegraphics[width=1\linewidth]{exact.128.256.512@500/exact128.png}\newline

    \begin{table}[H]
    \centering
    \small
    \setlength{\tabcolsep}{3pt}
    \begin{tabular}{l r r r r r}
    \hline
    \textbf{Group} & \textbf{Queries} & \textbf{Recall@5} & \textbf{Recall@10} & \textbf{mAP@10} & \textbf{Time(ms)} \\
    \hline
    all & 23 & 0.7304 & 0.7913 & 0.7441 & 3.9490 \\
    illumination & 11 & 0.6727 & 0.7636 & 0.7067 & 3.9058 \\
    viewpoint & 12 & 0.7833 & 0.8167 & 0.7783 & 3.9886 \\
    \hline
    \end{tabular}
    \par\smallskip
    \footnotesize Queries με λιγότερα από 10 αποτελέσματα: 0
    \end{table}
    
    % \includegraphics[width=1\linewidth]{exact.128.256.512@500/exact512.png}
    
    \item LSH K = 128 and K = 512:
    \begin{table}[H]
    \centering
    \small
    \setlength{\tabcolsep}{3pt}
    \begin{tabular}{l r r r r r}
    \hline
    \textbf{Group} & \textbf{Queries} & \textbf{Recall@5} & \textbf{Recall@10} & \textbf{mAP@10} & \textbf{Time(ms)} \\
    \hline
    all & 23 & 0.3565 & 0.3565 & 0.3493 & 0.0377 \\
    illumination & 11 & 0.2545 & 0.2545 & 0.2545 & 0.0378 \\
    viewpoint & 12 & 0.4500 & 0.4500 & 0.4361 & 0.0377 \\
    \hline
    \end{tabular}
    \par\smallskip
    \footnotesize Queries με λιγότερα από 10 αποτελέσματα: 1
    \end{table}
    
     % \includegraphics[width=1\linewidth]{lsh.128.256.512@500/lsh128.png}\newline

     \begin{table}[H]
    \centering
    \small
    \setlength{\tabcolsep}{3pt}
    \begin{tabular}{l r r r r r}
    \hline
    \textbf{Group} & \textbf{Queries} & \textbf{Recall@5} & \textbf{Recall@10} & \textbf{mAP@10} & \textbf{Time(ms)} \\
    \hline
    all & 23 & 0.3391 & 0.3391 & 0.3370 & 0.1727 \\
    illumination & 11 & 0.1818 & 0.1818 & 0.1818 & 0.1464 \\
    viewpoint & 12 & 0.4833 & 0.4833 & 0.4792 & 0.1967 \\
    \hline
    \end{tabular}
    \par\smallskip
    \footnotesize Queries με λιγότερα από 10 αποτελέσματα: 1
    \end{table}
     
      % \includegraphics[width=1\linewidth]{lsh.128.256.512@500/lsh512.png}
      
    \item Hypercube K = 128 and K = 512:
    \begin{table}[H]
    \centering
    \small
    \setlength{\tabcolsep}{3pt}
    kproj = 10, w = 0.05, M = 500, probes = 10
    \par\smallskip
    \begin{tabular}{l r r r r r}
    \hline
    \textbf{Group} & \textbf{Queries} & \textbf{Recall@5} & \textbf{Recall@10} & \textbf{mAP@10} & \textbf{Time(ms)} \\
    \hline
    all & 23 & 0.2696 & 0.2870 & 0.2636 & 0.0628 \\
    illumination & 11 & 0.2182 & 0.2182 & 0.2182 & 0.0655 \\
    viewpoint & 12 & 0.3167 & 0.3500 & 0.3051 & 0.0603 \\
    \hline
    \end{tabular}
    \par\smallskip
    \footnotesize Queries με λιγότερα από 10 αποτελέσματα: 0
    \end{table}
    
     % \includegraphics[width=1\linewidth]{hypercube.128.256.512@500/hypercube128.png}\newline

    \begin{table}[H]
    \centering
    \small
    \setlength{\tabcolsep}{3pt}
    kproj = 10, w = 0.05, M = 500, probes = 10
    \par\smallskip
    \begin{tabular}{l r r r r r}
    \hline
    \textbf{Group} & \textbf{Queries} & \textbf{Recall@5} & \textbf{Recall@10} & \textbf{mAP@10} & \textbf{Time(ms)} \\
    \hline
    all & 23 & 0.2522 & 0.2696 & 0.2559 & 0.1778 \\
    illumination & 11 & 0.1091 & 0.1091 & 0.1030 & 0.1504 \\
    viewpoint & 12 & 0.3833 & 0.4167 & 0.3961 & 0.2029 \\
    \hline
    \end{tabular}
    \par\smallskip
    \footnotesize Queries με λιγότερα από 10 αποτελέσματα: 0
    \end{table}
     
      % \includegraphics[width=1\linewidth]{hypercube.128.256.512@500/hypercube512.png}

    \item IVF FLAT K = 128 and K = 512:
    \begin{table}[H]
    \centering
    \small
    \setlength{\tabcolsep}{3pt}
    101 clusters, nprobe = 5
    \par\smallskip
    \begin{tabular}{l r r r r r}
    \hline
    \textbf{Group} & \textbf{Queries} & \textbf{Recall@5} & \textbf{Recall@10} & \textbf{mAP@10} & \textbf{Time(ms)} \\
    \hline
    all & 23 & 0.6696 & 0.6957 & 0.6783 & 0.3079 \\
    illumination & 11 & 0.6545 & 0.6909 & 0.6739 & 0.2807 \\
    viewpoint & 12 & 0.6833 & 0.7000 & 0.6824 & 0.3328 \\
    \hline
    \end{tabular}
    \par\smallskip
    \footnotesize Queries με λιγότερα από 10 αποτελέσματα: 0
    \end{table}
    
     % \includegraphics[width=1\linewidth]{ivf.128.256.512@500/ivf128.png}\newline

     \begin{table}[H]
    \centering
    \small
    \setlength{\tabcolsep}{3pt}
    101 clusters, nprobe = 5
    \par\smallskip
    \begin{tabular}{l r r r r r}
    \hline
    \textbf{Group} & \textbf{Queries} & \textbf{Recall@5} & \textbf{Recall@10} & \textbf{mAP@10} & \textbf{Time(ms)} \\
    \hline
    all & 23 & 0.7130 & 0.7565 & 0.7191 & 1.2189 \\
    illumination & 11 & 0.6364 & 0.6909 & 0.6545 & 1.1852 \\
    viewpoint & 12 & 0.7833 & 0.8167 & 0.7783 & 1.2497 \\
    \hline
    \end{tabular}
    \par\smallskip
    \footnotesize Queries με λιγότερα από 10 αποτελέσματα: 0
    \end{table}
     
     
      % \includegraphics[width=1\linewidth]{ivf.128.256.512@500/ivf512.png}

    \item IVF PQ K = 128 and K = 512:
    \begin{table}[H]
    \centering
    \small
    \setlength{\tabcolsep}{3pt}
    101 clusters, nprobe = 5, Μ = 16
    \par\smallskip
    \begin{tabular}{l r r r r r}
    \hline
    \textbf{Group} & \textbf{Queries} & \textbf{Recall@5} & \textbf{Recall@10} & \textbf{mAP@10} & \textbf{Time(ms)} \\
    \hline
    all & 23 & 0.6348 & 0.6870 & 0.6642 & 0.1366 \\
    illumination & 11 & 0.6364 & 0.6909 & 0.6617 & 0.1364 \\
    viewpoint & 12 & 0.6333 & 0.6833 & 0.6665 & 0.1367 \\
    \hline
    \end{tabular}
    \par\smallskip
    \footnotesize Queries με λιγότερα από 10 αποτελέσματα: 0
    \end{table}
    
     % \includegraphics[width=1\linewidth]{ivfpq.128.256.512@500/ivfpq128.png}\newline

     \begin{table}[H]
    \centering
    \small
    \setlength{\tabcolsep}{3pt}
    101 clusters, nprobe = 5, Μ = 16
    \par\smallskip
    \begin{tabular}{l r r r r r}
    \hline
    \textbf{Group} & \textbf{Queries} & \textbf{Recall@5} & \textbf{Recall@10} & \textbf{mAP@10} & \textbf{Time(ms)} \\
    \hline
    all & 23 & 0.6174 & 0.6870 & 0.6393 & 0.4838 \\
    illumination & 11 & 0.5091 & 0.5818 & 0.5331 & 0.4853 \\
    viewpoint & 12 & 0.7167 & 0.7833 & 0.7367 & 0.4825 \\
    \hline
    \end{tabular}
    \par\smallskip
    \footnotesize Queries με λιγότερα από 10 αποτελέσματα: 0
    \end{table}
     
      % \includegraphics[width=1\linewidth]{ivfpq.128.256.512@500/ivfpq512.png}
    
\end{itemize}

Παρατηρούμε πως, το Κ = 256 είναι η καλήτερη τιμή καθώς η ακρίβεια ανεβαίνει με αποδεκτό κόστος χρόνου, σε σχέση με το Κ = 128. Η αύξηση του περαιτέρω σε Κ = 512, οδηγεί στον μέσο χρόνο αναζήτησης να αυξάνεται χώρις όμως αυτό να έχει ανάλογα αποτελέσματα στην ακρίβεια. Σε συγκεκριμένες μετρήσεις, παρατηρούμε μάλιστα πως η ακρίβεια μένει στάστιμη ή και επηρεάζεται αρνητικά. 

\section{Επιρροή του ποσού των τυχαίων επιλεγμένων Sift Descriptors όταν αλλάζει σε S = 1000}

\begin{itemize}
    \item Exact Search:
    \begin{table}[H]
    \centering
    \small
    \setlength{\tabcolsep}{3pt}
    \begin{tabular}{l r r r r r}
    \hline
    \textbf{Group} & \textbf{Queries} & \textbf{Recall@5} & \textbf{Recall@10} & \textbf{mAP@10} & \textbf{Time(ms)} \\
    \hline
    all & 23 & 0.7478 & 0.8000 & 0.7597 & 1.9284 \\
    illumination & 11 & 0.6727 & 0.7636 & 0.6990 & 1.9141 \\
    viewpoint & 12 & 0.8167 & 0.8333 & 0.8154 & 1.9414 \\
    \hline
    \end{tabular}
    \par\smallskip
    \footnotesize Queries με λιγότερα από 10 αποτελέσματα: 0
    \end{table}

    % \includegraphics[width=1\linewidth]{exact.S1000/exact.s1000.png}\newline

    \item LSH:
    \begin{table}[H]
    \centering
    \small
    \setlength{\tabcolsep}{3pt}
    \begin{tabular}{l r r r r r}
    \hline
    \textbf{Group} & \textbf{Queries} & \textbf{Recall@5} & \textbf{Recall@10} & \textbf{mAP@10} & \textbf{Time(ms)} \\
    \hline
    all & 23 & 0.3130 & 0.3130 & 0.3101 & 0.0856 \\
    illumination & 11 & 0.1818 & 0.1818 & 0.1818 & 0.0787 \\
    viewpoint & 12 & 0.4333 & 0.4333 & 0.4278 & 0.0919 \\
    \hline
    \end{tabular}
    \par\smallskip
    \footnotesize Queries με λιγότερα από 10 αποτελέσματα: 0
    \end{table}

    % \includegraphics[width=1\linewidth]{lsh.S1000/lsh.s1000.png}\newline

    \item Hypercube: \newline
    \begin{table}[H]
    \centering
    \small
    \setlength{\tabcolsep}{3pt}
    kproj = 10, w = 0.05, M = 500, probes = 10
    \par\smallskip
    \begin{tabular}{l r r r r r}
    \hline
    \textbf{Group} & \textbf{Queries} & \textbf{Recall@5} & \textbf{Recall@10} & \textbf{mAP@10} & \textbf{Time(ms)} \\
    \hline
    all & 23 & 0.2609 & 0.2696 & 0.2617 & 0.1102 \\
    illumination & 11 & 0.1818 & 0.1818 & 0.1818 & 0.1088 \\
    viewpoint & 12 & 0.3333 & 0.3500 & 0.3350 & 0.1115 \\
    \hline
    \end{tabular}
    \par\smallskip
    \footnotesize Queries με λιγότερα από 10 αποτελέσματα: 0
    \end{table}

    % \includegraphics[width=1\linewidth]{hypercube.S1000/hypercube.s1000.png}\newline

    \item IVF FLAT:
    \begin{table}[H]
    \centering
    \small
    \setlength{\tabcolsep}{3pt}
    101 clusters, nprobe = 5
    \par\smallskip
    \begin{tabular}{l r r r r r}
    \hline
    \textbf{Group} & \textbf{Queries} & \textbf{Recall@5} & \textbf{Recall@10} & \textbf{mAP@10} & \textbf{Time(ms)} \\
    \hline
    all & 23 & 0.7130 & 0.7565 & 0.7271 & 0.4639 \\
    illumination & 11 & 0.6182 & 0.6909 & 0.6503 & 0.4577 \\
    viewpoint & 12 & 0.8000 & 0.8167 & 0.7975 & 0.4695 \\
    \hline
    \end{tabular}
    \par\smallskip
    \footnotesize Queries με λιγότερα από 10 αποτελέσματα: 0
    \end{table}

    % \includegraphics[width=1\linewidth]{ivf.S1000/ivf.s1000.png}\newline

    \item IVF PQ:
    \begin{table}[H]
    \centering
    \small
    \setlength{\tabcolsep}{3pt}
    101 clusters, nprobe = 5, M = 16
    \par\smallskip
    \begin{tabular}{l r r r r r}
    \hline
    \textbf{Group} & \textbf{Queries} & \textbf{Recall@5} & \textbf{Recall@10} & \textbf{mAP@10} & \textbf{Time(ms)} \\
    \hline
    all & 23 & 0.6783 & 0.7043 & 0.6852 & 0.2275 \\
    illumination & 11 & 0.6182 & 0.6545 & 0.6255 & 0.2271 \\
    viewpoint & 12 & 0.7333 & 0.7500 & 0.7400 & 0.2278 \\
    \hline
    \end{tabular}
    \par\smallskip
    \footnotesize Queries με λιγότερα από 10 αποτελέσματα: 0
    \end{table}

    % \includegraphics[width=1\linewidth]{ivfpq.S1000/ivfpq.s1000.png}\newline
    
\end{itemize}

Καθώς αυξάνεται ο αριθμός των τυχαίων Sift Descriptors, ο χρόνος αναζήτησης όπως περιμέναμε δεν επηρεάζεται, και υπάρχει βελτίωση σε όλους του αλγορίθμους στο Recall@n, mAP@10. Εξαίρεση είναι ο LSH, όπου όταν χρησημοποιούμε τις αρχικές τιμές: πλήθος συναρτήσεων ανά πίνακα = 5, πίνακες = 4, μέγεθος κουβά = 0.05, υπάρχει μείωση στην ακρίβεια του ενώ όσο μεγαλώνουμε τις τιμές αυτές, o LSH συμπεριφέρεται σαν του υπόλοιπους αλγορίθμους και παρουσιάζει βελτίωση.

\section{Αλλαγή αποδόσεων LSH}

Μεταβολή μεγέθους κελιού: -w
\begin{table}[H]
    \centering
    \footnotesize
    \setlength{\tabcolsep}{3pt}
    k = 4, L = 5, validation set, D = 10000
    \par\smallskip
    \begin{tabular}{r r r r r r r}
    \hline
    \textbf{w} & \textbf{Recall@5} & \textbf{Recall@10} & \textbf{mAP@10} & \textbf{mAP i\_} & \textbf{mAP v\_} & \textbf{Time(ms)} \\
    \hline
    0.05 & 0.3478 & 0.3478 & 0.3443 & 0.3564 & 0.3333 & 0.0496 \\
    0.10 & 0.6609 & 0.6870 & 0.6684 & 0.5673 & 0.7611 & 0.3989 \\
    0.15 & 0.7304 & 0.7391 & 0.7262 & 0.6382 & 0.8069 & 0.7815 \\
    0.20 & 0.7391 & 0.7565 & 0.7386 & 0.6641 & 0.8069 & 1.0725 \\
    0.30 & 0.7565 & 0.7739 & 0.7549 & 0.6800 & 0.8236 & 1.3565 \\
    \hline
    \end{tabular}
    \par\smallskip
    \footnotesize Queries με λιγότερα από 10 αποτελέσματα: 1 (w = 0.05), 0 στις υπόλοιπες
\end{table}

Παρατηρούμε ότι για μικρά w, τα buckets που δημιουργεί ο αλγόριθμος είναι πολύ μικρά και αυτό προκαλεί μεγάλη ανακρίβεια. Όσο αυξάνουμε το w η απόδοση του αλγόριθμου βελτιώνεται αισθητά και πλησιάζει όλο και περισσότερο την εξαντλητική αναζήτηση, αλλά αυξάνεται και ο χρόνος αναζήτησης. Βλέπουμε ότι για w=0.3 το LSH έχει το ίδιο mAP με την εξαντλητική αναζήτηση.

Μεταβολή πλήθους πινάκων: -L
\begin{table}[H]
    \centering
    \footnotesize
    \setlength{\tabcolsep}{3pt}
    k = 4, w = 0.15, validation set, D = 10000
    \par\smallskip
    \begin{tabular}{r r r r r r r}
    \hline
    \textbf{L} & \textbf{Recall@5} & \textbf{Recall@10} & \textbf{mAP@10} & \textbf{mAP i\_} & \textbf{mAP v\_} & \textbf{Time(ms)} \\
    \hline
    2  & 0.5304 & 0.5478 & 0.5327 & 0.4396 & 0.6181 & 0.3828 \\
    5  & 0.7304 & 0.7391 & 0.7262 & 0.6382 & 0.8069 & 0.7853 \\
    8  & 0.7391 & 0.7478 & 0.7327 & 0.6517 & 0.8069 & 0.9382 \\
    12 & 0.7478 & 0.7652 & 0.7468 & 0.6813 & 0.8069 & 1.2276 \\
    \hline
    \end{tabular}
    \par\smallskip
    \footnotesize Queries με λιγότερα από 10 αποτελέσματα: 0 σε όλες τις εκτελέσεις
\end{table}

Αντίστοιχα με τα w, και η αύξηση των πινάκων κατακερματισμού, αυξάνει την απόδοση του αλγορίθμου, ενώ επίσης παράλληλα αυξάνεται και ο χρόνος αναζήτησης.

\section{Αλλαγή αποδόσεων Hypercube}

Μεταβολή μέγιστου πλήθους υποψηφίων: -M
\begin{table}[H]
    \centering
    \footnotesize
    \setlength{\tabcolsep}{3pt}
    kproj = 12, w = 0.2, probes = 50, validation set, D = 10000
    \par\smallskip
    \begin{tabular}{r r r r r r r}
    \hline
    \textbf{M} & \textbf{Recall@5} & \textbf{Recall@10} & \textbf{mAP@10} & \textbf{mAP i\_} & \textbf{mAP v\_} & \textbf{Time(ms)} \\
    \hline
    500   & 0.1826 & 0.1826 & 0.1826 & 0.1455 & 0.2167 & 0.0993 \\
    1000  & 0.3217 & 0.3304 & 0.3239 & 0.2682 & 0.3750 & 0.3586 \\
    2000  & 0.5478 & 0.5565 & 0.5381 & 0.5176 & 0.5569 & 0.4033 \\
    5000  & 0.6957 & 0.7130 & 0.6938 & 0.6612 & 0.7236 & 0.9128 \\
    10000 & 0.7130 & 0.7217 & 0.7086 & 0.6558 & 0.7569 & 1.0980 \\
    \hline
    \end{tabular}
    \par\smallskip
    \footnotesize Queries με λιγότερα από 10 αποτελέσματα: 0 σε όλες τις εκτελέσεις
\end{table}

Παρατηρούμε ότι όσο αυξάνεται ο αριθμός των υποψήφιων σημείων που ελέγχει ο αλγόριθμος τόσο βελτιώνεται και η ακρίβεια του. Αυτό οφείλεται στο γεγονός ότι για μικρές τιμές του Μ, ο αλγόριθμος εξετάζει λίγες εικόνες από κάθε κορυφή που επισκέπτεται, με αποτέλεσμα να μην εξετάζονται πάντα όλες οι σχετικές εικόνες.

Μεταβολή πλήθους κορυφών που εξετάζονται: -probes
\begin{table}[H]
    \centering
    \footnotesize
    \setlength{\tabcolsep}{3pt}
    kproj = 12, w = 0.2, M = 10000, validation set, D = 10000
    \par\smallskip
    \begin{tabular}{r r r r r r r}
    \hline
    \textbf{probes} & \textbf{Recall@5} & \textbf{Recall@10} & \textbf{mAP@10} & \textbf{mAP i\_} & \textbf{mAP v\_} & \textbf{Time(ms)} \\
    \hline
    10  & 0.5478 & 0.5565 & 0.5454 & 0.4418 & 0.6403 & 0.5650 \\
    20  & 0.7043 & 0.7217 & 0.7025 & 0.6430 & 0.7569 & 1.2321 \\
    50  & 0.7130 & 0.7217 & 0.7086 & 0.6558 & 0.7569 & 1.2347 \\
    100 & 0.7391 & 0.7565 & 0.7385 & 0.6638 & 0.8069 & 1.3976 \\
    200 & 0.7478 & 0.7652 & 0.7476 & 0.6829 & 0.8069 & 1.3931 \\
    \hline
    \end{tabular}
    \par\smallskip
    \footnotesize Queries με λιγότερα από 10 αποτελέσματα: 0 σε όλες τις εκτελέσεις
\end{table}

Παρατηρούμε ότι η μεγαλύτερη αύξηση στην απόδοση του αλγορίθμου εμφανίζεται όταν το probes αλλάζει από 10 σε 20, ενώ μετά η αύξηση είναι μικρή. Παρόλα αυτά για μεγάλα probes (100 και 200) η απόδοση του αλγορίθμου φτάνει πολύ κοντά στην εξαντλητική αναζήτηση, αλλά χρειάζεται σχεδόν τον ίδιο χρόνο.

\section{Αλλαγή αποδόσεων IVFFlat}

Μεταβολή πλήθους λιστών που εξετάζονται: -nprobe
\begin{table}[H]
    \centering
    \footnotesize
    \setlength{\tabcolsep}{3pt}
    101 clusters, validation set, D = 10000, K = 256, S = 500
    \par\smallskip
    \begin{tabular}{r r r r r r r}
    \hline
    \textbf{nprobe} & \textbf{Recall@5} & \textbf{Recall@10} & \textbf{mAP@10} & \textbf{mAP i\_} & \textbf{mAP v\_} & \textbf{Time(ms)} \\
    \hline
    1   & 0.4435 & 0.4522 & 0.4461 & 0.3455 & 0.5383 & 0.1470 \\
    2   & 0.5826 & 0.5826 & 0.5765 & 0.4691 & 0.6750 & 0.2533 \\
    5   & 0.7304 & 0.7391 & 0.7256 & 0.6550 & 0.7903 & 0.4609 \\
    10  & 0.7478 & 0.7652 & 0.7468 & 0.6631 & 0.8236 & 0.7707 \\
    20  & 0.7565 & 0.7739 & 0.7549 & 0.6800 & 0.8236 & 1.2488 \\
    40  & 0.7565 & 0.7739 & 0.7549 & 0.6800 & 0.8236 & 1.7748 \\
    101 & 0.7565 & 0.7739 & 0.7549 & 0.6800 & 0.8236 & 1.9649 \\
    \hline
    \end{tabular}
    \par\smallskip
    \footnotesize Queries με λιγότερα από 10 αποτελέσματα: 0 σε όλες τις εκτελέσεις
    \end{table}

Παρατηρούμε ότι η αύξηση των clusters που ελέγχουμε φέρνει την μέθοδο όλο και πιο κοντά στα αποτελέσματα της εξαντλητικής μεθόδου και φτάνει την ίδια απόδοση για nprobe=20, όπου από εκεί και πέρα σταματά να βελτιώνεται και απλά η αναζήτηση απαιτεί περισσότερο χρόνο. Παρατηρούμε επίσης ότι για nprobe=10 είναι πολύ κοντά στην εξαντλητική χρειάζοντας μόνο 0.7707 ms.

Μεταβολή πλήθους λιστών: -kclusters
\begin{table}[H]
    \centering
    \footnotesize
    \setlength{\tabcolsep}{3pt}
    nprobe = 10, validation set, D = 10000, K = 256, S = 500
    \par\smallskip
    \begin{tabular}{r r r r r r r}
    \hline
    \textbf{kclusters} & \textbf{Recall@5} & \textbf{Recall@10} & \textbf{mAP@10} & \textbf{mAP i\_} & \textbf{mAP v\_} & \textbf{Time(ms)} \\
    \hline
    50  & 0.7565 & 0.7739 & 0.7549 & 0.6800 & 0.8236 & 1.2707 \\
    101 & 0.7478 & 0.7652 & 0.7468 & 0.6631 & 0.8236 & 0.8146 \\
    200 & 0.7478 & 0.7652 & 0.7468 & 0.6813 & 0.8069 & 0.4622 \\
    \hline
    \end{tabular}
    \par\smallskip
    \footnotesize Queries με λιγότερα από 10 αποτελέσματα: 0 σε όλες τις εκτελέσεις
    \end{table}

Παρατηρούμε ότι η αύξηση των clusters που χρησιμοποιούμε σε 200, μειώνει τον χρόνο αναζήτησης στο μισό (0.4622ms), ενώ συνεχίζει να έχει παρόμοια απόδοση με όταν τα clusters είναι 101. Παράλληλα, χρησιμοποιώντας 50 clusters, έχουμε παρόμοια απόδοση με την εξαντλητική αναζήτηση, αλλά με χρόνο 1.2707 ms.

\section{Αλλαγή αποδόσεων IVFPQ}

Μεταβολή πλήθους λιστών που εξετάζονται: -nprobe
\begin{table}[H]
    \centering
    \footnotesize
    \setlength{\tabcolsep}{3pt}
    101 clusters, M = 16, nbits = 8, validation set, D = 10000
    \par\smallskip
    \begin{tabular}{r r r r r r r}
    \hline
    \textbf{nprobe} & \textbf{Recall@5} & \textbf{Recall@10} & \textbf{mAP@10} & \textbf{mAP i\_} & \textbf{mAP v\_} & \textbf{Time(ms)} \\
    \hline
    1  & 0.4174 & 0.4348 & 0.4257 & 0.3455 & 0.4992 & 0.0629 \\
    2  & 0.5391 & 0.5565 & 0.5471 & 0.4691 & 0.6185 & 0.0885 \\
    5  & 0.6435 & 0.6783 & 0.6606 & 0.6200 & 0.6979 & 0.2008 \\
    10 & 0.6435 & 0.6783 & 0.6606 & 0.6200 & 0.6979 & 0.3298 \\
    20 & 0.6522 & 0.6870 & 0.6693 & 0.6382 & 0.6979 & 0.6294 \\
    \hline
    \end{tabular}
    \par\smallskip
    \footnotesize Queries με λιγότερα από 10 αποτελέσματα: 0 σε όλες τις εκτελέσεις
\end{table}

Παρατηρούμε ότι για τις ίδιες τιμές nprobe ο αλγόριθμος είναι λιγότερο ακριβής από τον IVFFlat, όμως σε πολύ λιγότερο χρόνο. Αυτό προκύπτει από τον περιορισμό που υπάρχει στο M (16) το οποίο αποδεικνύεται και από τα αποτελέσματα του παρακάτω πίνακα.

Μεταβολή πλήθους υποδιανυσμάτων: -M
\begin{table}[H]
    \centering
    \footnotesize
    \setlength{\tabcolsep}{3pt}
    101 clusters, nprobe = 10, nbits = 8, validation set, D = 10000
    \par\smallskip
    \begin{tabular}{r r r r r r r}
    \hline
    \textbf{M} & \textbf{Recall@5} & \textbf{Recall@10} & \textbf{mAP@10} & \textbf{mAP i\_} & \textbf{mAP v\_} & \textbf{Time(ms)} \\
    \hline
    8  & 0.5913 & 0.6435 & 0.6018 & 0.5397 & 0.6587 & 0.4220 \\
    16 & 0.6435 & 0.6783 & 0.6606 & 0.6200 & 0.6979 & 0.3298 \\
    32 & 0.6783 & 0.7652 & 0.7185 & 0.6699 & 0.7630 & 0.4510 \\
    64 & 0.7217 & 0.7652 & 0.7396 & 0.6679 & 0.8054 & 0.8016 \\
    \hline
    \end{tabular}
    \par\smallskip
    \footnotesize Queries με λιγότερα από 10 αποτελέσματα: 0 σε όλες τις εκτελέσεις
\end{table}

Παρατηρούμε ότι αυξάνοντας τον αριθμό των υποδιανυσμάτων η απόδοση αυξάνεται σταθερά, σε σημείο που για M=64,ο αλγόριθμος πλησιάζει πολύ την απόδοση του IVFFlat.


\section{Τελικά αποτελέσματα (test set)}

Χρησιμοποιούμε τις παραμέτρους που επιλέχθηκαν από τα πειράματα στο validation set (με D = 10000) με βάση την απόδοση τους. Χρησιμοποιήθηκε K = 256, S = 500, seed = 1. 

\begin{table}[H]
    \centering
    \footnotesize
    \setlength{\tabcolsep}{3pt}
    \begin{tabular}{l l}
    \hline
    \textbf{Μέθοδος} & \textbf{Παράμετροι} \\
    \hline
    LSH       & k = 4, L = 5, w = 0.15 \\
    Hypercube & kproj = 12, w = 0.2, M = 10000, probes = 20 \\
    IVFFlat   & kclusters = $\lceil\sqrt{n}\rceil$, nprobe = 10 \\
    IVFPQ  & kclusters = $\lceil\sqrt{n}\rceil$, nprobe = 10, M = 64, nbits = 8 \\
    \hline
    \end{tabular}
\end{table}

\begin{table}[H]
    \centering
    \footnotesize
    \setlength{\tabcolsep}{3pt}
    test set, D = 1000 (34 clusters)
    \par\smallskip
    \begin{tabular}{l r r r r r r}
    \hline
    \textbf{Method} & \textbf{Recall@5} & \textbf{Recall@10} & \textbf{mAP@10} & \textbf{mAP i\_} & \textbf{mAP v\_} & \textbf{Time(ms)} \\
    \hline
    Exact     & 0.7043 & 0.7826 & 0.7314 & 0.7085 & 0.7524 & 0.1681 \\
    LSH       & 0.6435 & 0.7130 & 0.6721 & 0.6236 & 0.7164 & 0.0874 \\
    Hypercube & 0.5217 & 0.5739 & 0.5328 & 0.4068 & 0.6482 & 0.0535 \\
    IVFFlat   & 0.7043 & 0.7826 & 0.7314 & 0.7085 & 0.7524 & 0.1551 \\
    IVFPQ  & 0.7130 & 0.7652 & 0.7201 & 0.6939 & 0.7441 & 0.4525 \\
    \hline
    \end{tabular}
    \par\smallskip
    \footnotesize Queries με λιγότερα από 10 αποτελέσματα: 0 σε όλες τις μεθόδους
\end{table}

\begin{table}[H]
    \centering
    \footnotesize
    \setlength{\tabcolsep}{3pt}
    test set, D = 5000 (72 clusters)
    \par\smallskip
    \begin{tabular}{l r r r r r r}
    \hline
    \textbf{Method} & \textbf{Recall@5} & \textbf{Recall@10} & \textbf{mAP@10} & \textbf{mAP i\_} & \textbf{mAP v\_} & \textbf{Time(ms)} \\
    \hline
    Exact     & 0.6870 & 0.7478 & 0.7056 & 0.7024 & 0.7085 & 0.7642 \\
    LSH       & 0.6435 & 0.7043 & 0.6662 & 0.6236 & 0.7051 & 0.3836 \\
    Hypercube & 0.5130 & 0.5391 & 0.5191 & 0.3955 & 0.6325 & 0.2386 \\
    IVFFlat   & 0.6783 & 0.7391 & 0.6969 & 0.6842 & 0.7085 & 0.5314 \\
    IVFPQ  & 0.6783 & 0.7217 & 0.6776 & 0.6736 & 0.6813 & 0.5595 \\
    \hline
    \end{tabular}
    \par\smallskip
    \footnotesize Queries με λιγότερα από 10 αποτελέσματα: 0 σε όλες τις μεθόδους
\end{table}

\begin{table}[H]
    \centering
    \footnotesize
    \setlength{\tabcolsep}{3pt}
    test set, D = 10000 (101 clusters)
    \par\smallskip
    \begin{tabular}{l r r r r r r}
    \hline
    \textbf{Method} & \textbf{Recall@5} & \textbf{Recall@10} & \textbf{mAP@10} & \textbf{mAP i\_} & \textbf{mAP v\_} & \textbf{Time(ms)} \\
    \hline
    Exact     & 0.6783 & 0.7391 & 0.6949 & 0.6933 & 0.6963 & 1.5912 \\
    LSH       & 0.6435 & 0.6870 & 0.6586 & 0.6145 & 0.6990 & 0.7517 \\
    Hypercube & 0.5130 & 0.5391 & 0.5165 & 0.3927 & 0.6300 & 0.5045 \\
    IVFFlat   & 0.6696 & 0.7304 & 0.6862 & 0.6752 & 0.6963 & 0.7338 \\
    IVFPQ  & 0.6609 & 0.7304 & 0.6784 & 0.6788 & 0.6780 & 0.6594 \\
    \hline
    \end{tabular}
    \par\smallskip
    \footnotesize Queries με λιγότερα από 10 αποτελέσματα: 0 σε όλες τις μεθόδους
\end{table}

\begin{table}[H]
    \centering
    \footnotesize
    \setlength{\tabcolsep}{3pt}
    test set, D = 25000 (24985 διαθέσιμες MIRFlickr, 159 clusters)
    \par\smallskip
    \begin{tabular}{l r r r r r r}
    \hline
    \textbf{Method} & \textbf{Recall@5} & \textbf{Recall@10} & \textbf{mAP@10} & \textbf{mAP i\_} & \textbf{mAP v\_} & \textbf{Time(ms)} \\
    \hline
    Exact     & 0.6609 & 0.7130 & 0.6811 & 0.6873 & 0.6754 & 3.8935 \\
    LSH       & 0.6261 & 0.6783 & 0.6445 & 0.6145 & 0.6719 & 2.1217 \\
    Hypercube & 0.5043 & 0.5391 & 0.5150 & 0.3927 & 0.6271 & 1.7111 \\
    IVFFlat   & 0.6609 & 0.7130 & 0.6811 & 0.6873 & 0.6754 & 1.3770 \\
    IVFPQ  & 0.6609 & 0.7130 & 0.6799 & 0.6909 & 0.6699 & 0.8554 \\
    \hline
    \end{tabular}
    \par\smallskip
    \footnotesize Queries με λιγότερα από 10 αποτελέσματα: 0 σε όλες τις μεθόδους
\end{table}

Με βάση τα αποτελέσματα του test set, παρατηρούμε ότι η εξαντλητική αναζήτηση συνεχίζει να έχει μεγάλη ακρίβεια, αλλά ο χρόνος αναζήτησης αυξάνεται σημαντικά όσο αυξάνεται και η βάση των εικόνων. Ο IVFFlat έχει επίσης αρκετά υψηλή ακρίβεια, σε σημείο που σε μερικές περιπτώσεις είναι ίση με της εξαντλητικής, ενώ για D = 25000 ο IVFPQ πετυχαίνει ακρίβεια σχεδόν ίση με της εξαντλητικής αναζήτησης, σε πολύ λιγότερο χρόνο. Παράλληλα, οι μεθόδοι Hypercube και LSH έχουν μικρούς χρόνους αναζήτησης αλλά είναι πολύ ανακριβής.



\end{document}
